\subsection{The characterization of separability by automorphisms}

THEOREM 4. --- Let \(\Omega\) be an algebraically closed extension of a field \(K\) and \(L\) a subextension of \(\Omega\) . Then the following conditions are equivalent:

\begin{enumerate}
\def\labelenumi{\alph{enumi})}
\item
  L is separable over K.
\item
  The intersection of the kernels of the \(\Omega\) -algebra homomorphisms of \(\Omega \otimes_{K} L\) into \(\Omega\) is reduced to 0.
\item
  For any elements \(a_1, \ldots, a_n\) of L linearly independent over K there exist K-automorphisms \(\sigma_1, \ldots, \sigma_n\) of \(\Omega\) such that \(\det(\sigma_i(a_j)) \neq 0\) .
\item
  Let V be a vector sub-K-space of L of finite dimension. Every K-linear mapping of V into \(\Omega\) is a linear combination (with coefficients in \(\Omega\) ) of the restrictions to V of K-automorphisms of \(\Omega\) .
\end{enumerate}

\(d) \Rightarrow c)\) : Let \(a_1, \ldots, a_n\) be elements of L linearly independent over K and let V be the vector sub-K-space of L generated by \(a_1, \ldots, a_n\) . The mapping \(f \mapsto (f(a_1), \ldots, f(a_n))\) is an \(\Omega\) -linear bijection of \(\mathrm{Hom}_{\mathbf{K}}(\mathbf{V}, \Omega)\) onto \(\Omega^n\) . Suppose that \(d)\) holds, then there exist K-automorphisms \(\sigma_1, \ldots, \sigma_n\) of \(\Omega\) such that the elements \((\sigma_i(a_1), \ldots, \sigma_i(a_n))\) of \(\Omega^n\) (for \(1 \leqslant i \leqslant n\) ) form a basis of \(\Omega^n\) . We thus have \(\det(\sigma_i(a_j)) \neq 0\) , so \(d)\) implies \(c)\) .

\(c) \Rightarrow b)\) : Assume that \(c)\) holds and let \(x\) be in \(\Omega \otimes_{\mathbb{K}} L\) . We write \(x\) in the form \(\sum_{j=1}^{n} x_j \otimes a_j\) with \(x_1, \ldots, x_n\) in \(\Omega\) and the elements \(a_1, \ldots, a_n\) of \(L\) linearly

independent over K. Choose K-automorphisms \(\sigma_{1},\ldots,\sigma_{n}\) of \(\Omega\) such that det \(\sigma_{i}(a_{j})\neq0\) ; let \(\chi_{1},\ldots,\chi_{n}\) be the \(\Omega\) -homomorphisms of \(\Omega\otimes_{K}L\) into \(\Omega\) such that \(\chi_{i}(a\otimes b)=a\cdot\sigma_{i}(b)\) for \(a\in\Omega\) and \(b\in L\) . Suppose that \(\chi_{i}(x)=0\) for \(1\leqslant i\leqslant n\) , in other words, that \(\sum_{j=1}^{n}x_{j}\cdot\sigma_{i}(a_{j})=0\) for \(1\leqslant i\leqslant n\) . Since we have assumed that the matrix \((\sigma_{i}(a_{j}))\) has a non-zero determinant, we have \(x_{i}=0\) for \(1\leqslant i\leqslant n\) and so \(x=0\) .

\(b) \Rightarrow a)\) : Since every extension of a field of characteristic 0 is separable (V, p.~122, Th. 1) it suffices to examine the case where K is of characteristic \(p \neq 0\) . Let X be the set of all \(\Omega\) -algebra homomorphisms of \(\Omega \otimes_{K} L\) into \(\Omega\) and \(f\) the homomorphism of \(\Omega \otimes_{K} L\) into \(\Omega^{X}\) defined by \(f(u) = (\chi(u))_{\chi \in X}\) for \(u \in \Omega \otimes_{K} L\) . Condition \(b)\) means that \(f\) is injective and so implies that the ring \(\Omega \otimes_{K} L\) is reduced. Hence condition \(b)\) of Th. 2 (V, p.~122) is satisfied with \(K' = \Omega\) , so L is separable over K.

\(a) \Rightarrow d)\) : Suppose that L is separable over K. Let V be a vector sub-K-space of finite dimension of L, \(V_{(\Omega)} = \Omega \otimes_K V\) the vector \(\Omega\) -space derived from V by extension of scalars and \(f_0\) the linear form on \(V_{(\Omega)}\) such that \(f_0(x \otimes y) = xy\) for \(x \in \Omega\) , \(y \in V\) . Denote by G the group of K-automorphisms of \(\Omega\) ; for \(\sigma \in G\) we put \(\sigma_V = \sigma \otimes Id_V\) and \(g_\sigma = \sigma \circ f_0 \circ \sigma_V^{-1}\) .

For each \(\sigma \in G\) the mapping \(g_{\sigma}\) of \(\mathbf{V}_{(\Omega)}\) into \(\Omega\) is \(\Omega\) -linear and maps \(x \otimes y\) to \(x \cdot \sigma(y)\) for \(x \in \Omega\) , \(y \in V\) . The kernel \(\mathbf{N}_{\sigma}\) of \(g_{\sigma}\) is therefore a vector subspace of \(\mathbf{V}_{(\Omega)}\) , and the same holds for \(\mathbf{N} = \bigcap_{\sigma \in G} \mathbf{N}_{\sigma}\) . If \(p\) is the characteristic

exponent of K, the field of invariants of G in \(\Omega\) is equal to \(\mathbf{K}^{p^{-\infty}}\) (V, p.~112, Prop. 10). We clearly have \(\sigma_{\mathsf{V}}(\mathbf{N}) = \mathbf{N}\) for all \(\sigma \in \mathbf{G}\) ; therefore (V, p.~63, Cor.) the vector \(\Omega\) -space N is generated by \(\mathbf{N}_0 = \mathbf{N} \cap (\mathbf{K}^{p^{-\infty}} \otimes \mathbf{V})\) . Since L is separable over K, the fields \(\mathbf{K}^{p^{-\infty}}\) and L are linearly disjoint over K (V, p.~123, Cor. 1); we have \(\mathbf{K}^{p^{-\infty}} \otimes_{\mathbb{K}} \mathbf{V} \subset \mathbf{K}^{p^{-\infty}} \otimes_{\mathbb{K}} \mathbf{L}\) and \(f_0(x \otimes y) = xy\) for \(x \in \Omega\) and \(y \in \mathbf{V}\) . It follows that the restriction of \(f_0\) to \(\mathbf{K}^{p^{-\infty}} \otimes_{\mathbb{K}} \mathbf{V}\) is injective. Now \(f_0 = g_1\) is zero on N and a fortiori on \(\mathbf{N}_0 \subset \mathbf{K}^{p^{-\infty}} \otimes_{\mathbb{K}} \mathbf{V}\) . We thus have \(\mathbf{N}_0 = 0\) , whence \(\mathbf{N} = 0\) . Since V is of finite dimension over K, \(\mathbf{V}_{(\Omega)}\) is of finite dimension over \(\Omega\) ; the intersection of the kernels of the linear forms \(g_\sigma\) is zero, hence (II, p.~301, Th. 7) the family \((g_\sigma)_{\sigma \in G}\) generates the dual of \(\mathbf{V}_{(\Omega)}\) . Now let \(u\) be a K-linear mapping of V into \(\Omega\) ; let \(\tilde{u}\) be the linear form on \(\mathbf{V}_{(\Omega)}\) which maps \(x \otimes y\) to \(x u(y)\) for \(x \in \Omega\) and \(y \in V\) . By what has been said there exist elements \(\sigma_1, ..., \sigma_n\) of G and \(\lambda_1, ..., \lambda_n\) of \(\Omega\) such that \(\tilde{u} = \sum_{i=1}^{n} \lambda_i g_{\sigma_i}\) , whence \(u(y) = \sum_{i=1}^{n} \lambda_i \sigma_i(y)\) for all \(y \in V\) . So we have shown that \(a)\) implies \(d)\) .

